\chapter*{Introduction}
\addcontentsline{toc}{chapter}{Introduction}
\markboth{Introduction}{Introduction}

\section*{Context and motivation}

Signalized intersections are among the most demanding places for an automated vehicle. A single
junction can hold a dozen traffic lights: circular lamps for the main flow, arrow lamps for turning
lanes, lights for the opposite approach, lights of the next intersection further down the road and
pedestrian signals. A human driver resolves this scene almost instantly by knowing which lane they are
in, where they intend to go, and which lights are meant for them. For a perception system, the same
reasoning has to be made explicit.

Most traffic light recognition systems answer a narrower question: \emph{where} are the traffic
lights, and \emph{what color} is each one. Detecting and classifying lights is necessary, but it is not
what the planner consumes. The planner needs to know whether the maneuver it intends to perform is
permitted \emph{now}. Between the two questions lies a relevance and association problem---which of
the detected lights govern the ego vehicle, and which maneuver each of them controls---that has
traditionally been solved with prior high-definition (HD) maps storing the position of every light and
the lanes it serves~\cite{fairfield2011traffic,possatti2019traffic,polley2025atlas}. Maps are accurate
where they exist, but they are expensive to build and maintain and they fail silently when the world
changes. Map-free alternatives infer relevance from the image, using annotated lane
metadata~\cite{langenberg2019deep}, detected lane arrows~\cite{polley2024tldready} or a joint
detection and image-level classification model~\cite{trinci2026color}; they decide which lights are
relevant, but still not which maneuver each of them permits.

A recent line of work removes the map and asks the network to answer the planner's question directly.
Zou et al.~\cite{zou2026oneperception} formulate \emph{maneuver-level signal intent prediction}: from
one front-camera image, predict, for each of the four maneuvers straight, left, right and U-turn, the
color of the signal that governs it and whether the maneuver may proceed. (This thesis drops the U-turn,
which almost always follows the left-turn signal; \cref{sec:prob-def}.) Their direction-aware model
uses four learnable \emph{direction queries} that attend to a global image token and to region tokens
from an auxiliary traffic-light detector, and it outperforms both structured detection pipelines and
multimodal large language models on a benchmark derived from OpenLane-V2~\cite{wang2023openlanev2}.

\section*{Research problem}

The proposal is attractive but leaves several questions open. The benchmark labels and the code are not
released, so the reported numbers cannot be checked directly and the label semantics---which light
governs a maneuver, how a turn without its own arrow is treated---have to be reconstructed from the
paper. It is also unclear whether the end-to-end direction-query design is needed at all: the structured
pipelines it was compared with were designed for a different output, and a well-tuned modern detector
combined with a simple, explicit association rule might do as well. Finally, the evaluation is limited to OpenLane-V2,
which mixes US and Singapore driving under one set of rules. Traffic signals look different across
countries, and the rules that link a lamp to a maneuver differ too: in Germany a right turn obeys the
circular signal, while in the United States a right turn on red is generally permitted unless a
dedicated signal says otherwise~\cite{fhwa2009mutcd,stvo}. How a model trained in one setting behaves in
another is the question that matters most for deployment.

\section*{Objectives and research questions}

The thesis pursues three objectives:
\begin{enumerate}
  \item to \textbf{rebuild} the maneuver-level benchmark of Zou et al.\ from public data only,
        documenting every assumption that the reconstruction requires;
  \item to build and evaluate a \textbf{structured model} from a state-of-the-art real-time detector and
        the same label rules, under a statistically sound protocol;
  \item to \textbf{measure robustness} across datasets, countries, traffic rules, illumination and
        cameras.
\end{enumerate}
These objectives translate into three research questions:
\begin{description}[leftmargin=1.2cm,style=nextline]
  \item[RQ1] Can the benchmark labels be rebuilt from the public OpenLane-V2 annotations, and does the
             result match the published label distribution?
  \item[RQ2] How well does a structured pipeline, in which a detector predicts light classes that encode
             the controlled maneuvers and an explicit rule assigns them, solve the task, and where does
             it fail?
  \item[RQ3] How well does a model trained on one dataset transfer to others, and which conditions
             cause the largest degradation?
\end{description}

\section*{Methodology in brief}

The work is empirical. Per-frame, per-maneuver labels are derived from the lane graph and the
lane--traffic-element links of OpenLane-V2, for both its Argoverse~2 (subset A) and nuScenes (subset
B) parts, and the derivation is validated against the label distribution published by Zou et al.\
and by manual spot-checks. Compatible label builders are written for the DriveU Traffic Light Dataset
(DTLD)~\cite{fregin2018dtld} and the LISA Traffic Light Dataset~\cite{jensen2016vision}. Models are
trained on a single consumer GPU; all hyperparameters are tuned on a held-out part of the training
scenes so that the evaluation split is never used for model selection. Metrics follow the original
paper (color accuracy, passability accuracy, exact match, their mean TSS, and the False-Go and
False-Stop rates), and are complemented with per-direction analyses and a segment-level paired
bootstrap for comparisons between models.

\section*{Contributions}

\begin{itemize}
  \item An open reconstruction of the maneuver-level benchmark from OpenLane-V2, with documented rules
        whose output reproduces the published label distribution, and evidence that the benchmark
        combines subsets A and B under US driving rules.
  \item A structured RF-DETR model whose classes encode the maneuvers each light controls, so that a
        perfect detector reproduces every label, reaching a TSS of 0.899 on subset A and 0.915 on
        subset B of the rebuilt benchmark.
  \item Label builders for DTLD and LISA that follow the same conventions, a model trained on DTLD and
        an out-of-distribution evaluation on LISA, OpenLane-V2 and ATLAS.
  \item A per-direction error analysis showing where aggregate scores hide safety-relevant failures.
\end{itemize}

\section*{Structure of the thesis}

\Cref{ch:background} reviews traffic light recognition, the map-based, map-free and topology-based
approaches to deciding which lights are relevant, the maneuver-level formulation, the detection and
detection architectures used later, and the public traffic-light datasets. \Cref{ch:problem} analyses the maneuver-level task in depth: its formal definition,
the traffic rules behind the labels, the metrics and their blind spots, and the two design paradigms
for solving it. \Cref{ch:method} describes the data, the label construction, the model, the training
and evaluation protocol, and presents all results. \Cref{ch:discussion} discusses the findings against
the original paper and the literature, the limitations of the work and directions for future research.
The Conclusions summarize the answers to the research questions. The appendices collect the label
rules in detail, the full hyperparameters, per-epoch results and the commands needed to reproduce every
number.
